\section*{Реферат}

Настоящий отчёт содержит \ztotpages\ страниц, \totalfigures\ русунков, \totaltables\ таблицы и 4 приложения.

Отчёт посвящён развёртыванию вычислительного кластера под управлением SLURM и разработке программы генерации пользовательских заданий на основе статистики реальной очереди. В ходе работы была подготовлена среда из пяти виртуальных машин на базе openSUSE Leap: управляющий узел и четыре вычислительных узла. Для узлов настроены статическая адресация, разрешение имён, SSH-доступ по ключам, общая служба аутентификации MUNGE и конфигурация SLURM.

Для генерации потока заданий использован фрагмент набора данных SLURM, отфильтрованный по индивидуальному варианту: UID 50728 и имя задания \mintinline{text}{qe7}. По исходным данным рассчитаны распределения длительности выполнения и ошибки пользовательского прогноза времени. На основе полученных параметров реализован генератор пользовательских заявок SLURM, который формирует задания с командой \mintinline{text}{sleep}, задаёт плановое время выполнения и автоматически помещает задания в очередь. Результаты представлены через статистику выполненных заданий, графики исходных и сгенерированных распределений, а также сравнение планового и фактического времени выполнения.

Ключевые слова: SLURM, MUNGE, openSUSE Leap, кластер, виртуальная машина, пользовательская заявка, очередь заданий, логнормальное распределение, нормальное распределение, генерация заданий.

\newpage
